\documentclass[11pt,a4paper]{article}

% ----- Encoding & fonts (Calibri-like sans-serif look) -----
\usepackage[utf8]{inputenc}
\usepackage[T1]{fontenc}
\usepackage{helvet}
\renewcommand{\familydefault}{\sfdefault}

% ----- Layout -----
\usepackage[margin=1in]{geometry}
\usepackage{array}
\usepackage{tabularx}
\usepackage{enumitem}
\usepackage{graphicx}
\usepackage{xcolor}
\usepackage{parskip}
\usepackage[hidelinks]{hyperref}

% ----- Helpers -----
% A blank-cell strut that gives every status-table row a fillable height.
\newcommand{\fillrow}{\rule[-0.9em]{0pt}{2.4em}}
% A short underline for fill-in fields (name / ID).
\newcommand{\fillin}[1]{\underline{\hspace{#1}}}
% An answer area: a label followed by reserved vertical space.
\newcommand{\answerspace}[1]{%
  \par\vspace{1mm}\noindent\textit{Answer / images:}\par\vspace{#1}%
}

\setlength{\parindent}{0pt}

\begin{document}

% ======================================================================
%  PART 0 - Title block (exact replica of the .docx template)
% ======================================================================
\begin{center}
  {\fontsize{22}{26}\selectfont \textbf{FPGA experiment report}}
\end{center}

\vspace{0.5em}
Task: FPGA 4

\vspace{1em}
Student 1\hspace{1.5em}name: \fillin{4cm}\hspace{1.5em}ID: \fillin{3cm}

Student 2:\hspace{1.5em}name: \fillin{4cm}\hspace{1.5em}ID: \fillin{3cm}

\vspace{1.5em}

% ======================================================================
%  PART 1 - Status update
% ======================================================================
{\fontsize{14}{17}\selectfont \textbf{Part 1 - Status update}}

\vspace{0.5em}
Fill in the following table according to your progress at the time of the submission.
Each row should refer to one of the modules together with its simulation.

The status column should be filled with one of the following:
\begin{itemize}[noitemsep,topsep=2pt]
  \item Done - If both are written and simulation passed.
  \item Partially done - If you started writing the module and simulation but it's not working properly.
  \item Not started - If you didn't get to that part of the task.
\end{itemize}

When the status is partially done, please use the notes/comments column to explain in a few more words what is the status of the module and what is the status of the simulation.

\vspace{0.5em}
Fill here:

\vspace{0.5em}
\begin{tabularx}{\textwidth}{|p{3.8cm}|p{3cm}|X|}
\hline
Module name/Step & status & notes/comments \\ \hline
\fillrow Control                       & & \\ \hline
\fillrow Debouncer                     & & \\ \hline
\fillrow Stopwatch                     & & \\ \hline
\fillrow Elaborated Design             & & \\ \hline
\fillrow Synthesis                     & & \\ \hline
\fillrow Setting up constraints        & & \\ \hline
\fillrow Implementing the design       & & \\ \hline
\fillrow FPGA programming and debugging & & \\ \hline
\end{tabularx}

\vspace{1.5em}
\textbf{Issues:}

In this section please describe the current issues you are facing.
\begin{itemize}[noitemsep,topsep=2pt]
  \item Per each module and simulation, if you receive any errors at the time of submission, please describe which errors and what they mean.
\end{itemize}

Fill here:
\begin{itemize}[topsep=2pt]
  \item Control:
  \item Debouncer:
  \item Stopwatch:
  \item Elaborated Design:
  \item Synthesis:
  \item Setting up constraints:
  \item Implementing the design:
  \item FPGA programming and debugging:
\end{itemize}

\vspace{1.5em}

% ======================================================================
%  PART 2 - Simulations and answers
% ======================================================================
{\fontsize{14}{17}\selectfont \textbf{Part 2 - Simulations and answers}}

\vspace{0.5em}
In this section, please add the images for each simulation as described in the report. Remember to use annotations and explain in your own words what is shown in each simulation.

In some of the tasks in the report, there are also questions that you should answer here.

Make sure to include images of all simulations running at the time of submission, even if they are not correct.

\vspace{0.5em}
Fill here:
\begin{itemize}[topsep=2pt]
  \item Control:
  \item Debouncer:
  \item Stopwatch:
  \item Elaborated Design:
  \item Synthesis:
  \item Setting up constraints:
  \item Implementing the design:
  \item FPGA programming and debugging:
\end{itemize}

\clearpage

% ======================================================================
%  PART 3 - Questions to fill and answer (from Experiment 4 assignment)
% ======================================================================
{\fontsize{14}{17}\selectfont \textbf{Part 3 - Questions to fill and answer}}

\vspace{0.3em}
The following are all the items from the \emph{Experiment FPGA 4 -- Dual Stopwatch} assignment that explicitly ask for content in the progress report. Answer each one below (add waveforms, screenshots and explanations where requested).

\vspace{0.8em}
\textbf{Task 2 -- Control}
\begin{enumerate}[label=(\roman*),leftmargin=2.2em,topsep=2pt]
  \item Provide an annotated waveform of your simulation. In addition, take a screenshot of the Tcl console window showing the ``Test Passed'' message and the ``\$finish called\ldots'' message. Explain what is shown in the simulation.
  \answerspace{3.5cm}
  \item The FSM is presented as a Mealy state diagram. Explain whether this FSM is a Moore FSM.
  \answerspace{3cm}
\end{enumerate}

\vspace{0.8em}
\textbf{Task 3 -- Debouncer}
\begin{enumerate}[label=(\alph*),leftmargin=2.2em,topsep=2pt,start=3]
  \item Discuss the tradeoff between choosing high and low values for the counter, and how this is related to the low-pass-filter debouncing approach we saw in the previous experiments.
  \answerspace{3.5cm}
\end{enumerate}

\vspace{0.8em}
\textbf{Task 4 -- Stopwatch}
\begin{enumerate}[label=(\alph*),leftmargin=2.2em,topsep=2pt]
  \item Design a block diagram for the Stopwatch module. Use the previous designs as instantiated modules (no need to show the internal logic); show and explain any other logic and connectivity.
  \answerspace{4cm}
\end{enumerate}

\vspace{0.8em}
\textbf{Task 5 -- Elaborated Design}
\begin{enumerate}[label=(\alph*),leftmargin=2.2em,topsep=2pt,start=2]
  \item Provide the screenshot of the elaborated design schematic.
  \answerspace{4cm}
\end{enumerate}

\vspace{0.8em}
\textbf{Task 6 -- Synthesis}
\begin{enumerate}[label=(\roman*),leftmargin=2.2em,topsep=2pt]
  \item Provide a screenshot of the synthesized design schematic.
  \answerspace{4cm}
  \item Describe the differences that appeared in the synthesized design compared to the elaborated design.
  \answerspace{3cm}
  \item Describe the issues you encountered during synthesis and how you solved them (errors, critical warnings, and warnings).
  \answerspace{3cm}
\end{enumerate}

\vspace{0.8em}
\textbf{Task 8 -- Implementing the design}
\begin{enumerate}[label=(\roman*),leftmargin=2.2em,topsep=2pt]
  \item Provide a screenshot of the design timing summary where you must exhibit positive slack both for setup and for hold times.
  \answerspace{4cm}
  \item Specify the lengths (in nanoseconds) of your critical setup path (i.e. the \emph{path length} with the lowest setup slack -- do not just specify the slack value).
  \answerspace{3cm}
\end{enumerate}

\vspace{0.8em}
\textbf{Task 9 -- FPGA programming and debugging}
\begin{enumerate}[label=(\alph*),leftmargin=2.2em,topsep=2pt,start=3]
  \item After your design was successfully approved, remove the debug core, re-implement the design, and provide a screenshot of your Project Manager $\rightarrow$ Project Summary screen including post-implementation Timing, Power, and Utilization details (table view).
  \answerspace{4cm}
\end{enumerate}

\end{document}
